\documentclass[10pt,a4paper]{article}

\usepackage[spanish,es-nodecimaldot]{babel}
\usepackage[utf8]{inputenc}
\usepackage[T1]{fontenc}
\usepackage{newtxtext,newtxmath}
\usepackage{geometry}
\usepackage{microtype}
\usepackage{amsmath,mathtools}
\usepackage{enumitem}
\usepackage{xcolor}
\usepackage{hyperref}
\usepackage{fancyhdr}
\usepackage[most]{tcolorbox}

\geometry{margin=2.5cm,includeheadfoot}
\setlength{\headheight}{15pt}
\setlength{\headsep}{0.28cm}
\setlength{\footskip}{0.62cm}
\setlength{\parindent}{0pt}
\setlength{\parskip}{0.28em}
\setlength{\abovedisplayskip}{0.5em}
\setlength{\belowdisplayskip}{0.5em}
\setlength{\abovedisplayshortskip}{0.25em}
\setlength{\belowdisplayshortskip}{0.25em}
\setlist[enumerate]{leftmargin=2.2em,itemsep=0.3em,topsep=0.25em}
\setlist[itemize]{leftmargin=1.5em,itemsep=0.2em,topsep=0.2em}

\definecolor{repblue}{gray}{0}
\definecolor{repteal}{gray}{0}
\definecolor{repgray}{gray}{0}
\definecolor{repaccent}{gray}{0}
\definecolor{repline}{gray}{0}
\definecolor{replight}{gray}{1}

\hypersetup{
  hidelinks,
  pdftitle={TRGF Padilla Tarea 4},
  pdfauthor={Juan Ignacio Padilla Barrientos}
}

\pagestyle{fancy}
\fancyhf{}
\fancyhead[L]{\textcolor{repblue}{\small\itshape TRGF Padilla}}
\fancyhead[R]{\textcolor{repgray}{\small Seminario UCR I 2026}}
\fancyfoot[C]{\textcolor{repgray}{\small \thepage}}
\renewcommand{\headrulewidth}{0.3pt}
\renewcommand{\footrulewidth}{0pt}

\newcommand{\CC}{\mathbb{C}}

\tcbset{
  enhanced,
  sharp corners,
  boxrule=0.45pt,
  arc=0pt,
  left=0.85em,
  right=0.85em,
  top=0.45em,
  bottom=0.4em,
  boxsep=0pt
}

\newcommand{\inlinehead}[2][repblue]{%
  \par\smallskip
  {\color{#1}\bfseries #2.}\hspace{0.45em}%
}

\newcommand{\topbanner}{%
  \begin{tcolorbox}[
    colback=replight,
    colframe=repline,
    borderline west={1.2pt}{0pt}{repblue},
    borderline east={1.2pt}{0pt}{repteal},
    left=1.0em,
    right=1.0em,
    top=0.7em,
    bottom=0.55em
  ]
    {\Large\bfseries\color{repblue} TRGF Padilla \textemdash\ Tarea 4}\\[-0.1em]
    {\normalsize\color{repteal} Seminario UCR I 2026}
  \end{tcolorbox}
}

\newcounter{probcnt}

\newcommand{\problem}[2]{%
  \stepcounter{probcnt}%
  \ifnum\value{probcnt}>1\newpage\fi
  \par\vspace{0.25em}
  {\large\bfseries\color{repblue} Problema #1}\par
  {\color{repline}\rule{\linewidth}{0.55pt}}\par
  \begin{tcolorbox}[
    breakable,
    colback=white,
    colframe=repline,
    borderline west={1.6pt}{0pt}{repteal},
    left=0.7em,
    right=0.7em,
    top=0.35em,
    bottom=0.3em,
    before skip=0.35em,
    after skip=0.45em
  ]
    #2
  \end{tcolorbox}
}

\newcommand{\workspace}[1]{%
  \par\vspace*{#1\baselineskip}
}

\begin{document}

\color{black}
\topbanner
\small

\problem{1}{%
Sea $g\in G$ cualquier elemento. Sea $Z_g$ el conjunto de elementos en $G$
que conmutan con $g$, y sea $C_g$ la clase conjugada de $g$. Demuestre que
\[
  |C_g|\,|Z_g| = |G|.
\]
En particular, el tamaño de $Z_g$ es el mismo para cualesquiera dos elementos
en la misma clase conjugada.
}

\inlinehead{Respuesta} Consideremos la aplicación
\[
  \phi\colon G \to C_g,
  \qquad
  \phi(x)=xgx^{-1}.
\]
Por definición de clase conjugada, $\phi$ es sobreyectiva. Además, para
$x,y\in G$ se tiene
\begin{align*}
  \phi(x)=\phi(y)
  &\iff xgx^{-1}=ygy^{-1} \\
  &\iff y^{-1}x\,g=g\,y^{-1}x \\
  &\iff y^{-1}x\in Z_g \\
  &\iff xZ_g=yZ_g.
\end{align*}
Así, dos elementos de $G$ tienen la misma imagen por $\phi$ exactamente cuando
pertenecen a la misma clase lateral izquierda de $Z_g$. Por lo tanto,
$\phi$ induce una biyección
\[
  G/Z_g \cong C_g.
\]
Como $G$ es finito, concluimos que
\[
  |C_g|=[G:Z_g]=\frac{|G|}{|Z_g|},
\]
es decir,
\[
  |C_g|\,|Z_g|=|G|.
\]

Para la última afirmación, si $h\in C_g$, existe $a\in G$ tal que
$h=aga^{-1}$. Entonces, para todo $z\in G$,
\begin{align*}
  z\in Z_g
  &\iff zg=gz \\
  &\iff (aza^{-1})(aga^{-1})=(aga^{-1})(aza^{-1}) \\
  &\iff aza^{-1}\in Z_h.
\end{align*}
De aquí se sigue que la conjugación por $a$ define una biyección
$Z_g \to Z_h$, $z\mapsto aza^{-1}$. En particular,
\[
  |Z_g|=|Z_h|
\]
si $g$ y $h$ están en la misma clase conjugada.\hfill$\blacksquare$

\problem{2}{%
Sea $(U,\pi)$ una representación irreducible de $G$. Sea $a\in G$ un elemento
central en $G$ (es decir, tal que $ag = ga$ para todo $g\in G$). Demuestre
que $\pi(a)$ debe ser un operador escalar, es decir, de la forma $\lambda I$.
}

\inlinehead{Respuesta} Sea $T:=\pi(a)$. Como $a$ es central en $G$, para todo
$g\in G$ se tiene
\[
  T\pi(g)=\pi(a)\pi(g)=\pi(ag)=\pi(ga)=\pi(g)\pi(a)=\pi(g)T.
\]
Es decir, $T$ conmuta con todos los operadores $\pi(g)$.

Como $U$ es un espacio vectorial complejo de dimensión finita, $T$ tiene al
menos un valor propio $\lambda\in\CC$. Sea
\[
  W:=\ker(T-\lambda I).
\]
Entonces $W\neq\{0\}$. Veamos que $W$ es $G$-invariante. Si $v\in W$, entonces
$Tv=\lambda v$, y para todo $g\in G$,
\[
  (T-\lambda I)\bigl(\pi(g)v\bigr)
  = T\pi(g)v-\lambda\pi(g)v
  = \pi(g)Tv-\lambda\pi(g)v
  = \pi(g)(Tv-\lambda v)=0.
\]
Por lo tanto, $\pi(g)v\in W$, y así $W$ es un subespacio $G$-invariante de
$U$.

Como $(U,\pi)$ es irreducible y $W\neq\{0\}$, se sigue que $W=U$. Esto
significa que
\[
  (T-\lambda I)u=0
  \qquad\text{para todo } u\in U,
\]
es decir,
\[
  \pi(a)=T=\lambda I.
\]
Concluimos que $\pi(a)$ es un operador escalar.\hfill$\blacksquare$

\problem{3}{%
Sean $(U,\pi)$ y $(V,\rho)$ dos representaciones de dimensión finita,
\textit{isomorfas}, de un grupo finito $G$. Demuestre que
\[
  \chi_U = \chi_V,
\]
es decir, que tienen el mismo carácter.
}

\inlinehead{Respuesta} Como $(U,\pi)$ y $(V,\rho)$ son isomorfas, existe un
isomorfismo lineal $T\colon U\to V$ tal que
\[
  T\pi(g)=\rho(g)T
  \qquad\text{para todo } g\in G.
\]
Como $T$ es invertible, esto equivale a
\[
  \rho(g)=T\pi(g)T^{-1}
  \qquad\text{para todo } g\in G.
\]
Por definición del carácter de una representación de dimensión finita,
\[
  \chi_U(g)=\operatorname{tr}(\pi(g))
  \qquad\text{y}\qquad
  \chi_V(g)=\operatorname{tr}(\rho(g)).
\]
Entonces, para todo $g\in G$,
\[
  \chi_V(g)=\operatorname{tr}(\rho(g))
  =\operatorname{tr}\bigl(T\pi(g)T^{-1}\bigr)
  =\operatorname{tr}(\pi(g))
  =\chi_U(g),
\]
donde usamos que la traza es invariante bajo conjugación. Por lo tanto,
$\chi_U=\chi_V$.\hfill$\blacksquare$

\problem{4}{%
Sea $G$ un grupo finito y sea $A=\CC[G]$. Suponga que $A$ se puede descomponer
como
\[
  A = U \oplus V,
\]
donde $U$ y $V$ son dos subrepresentaciones. Demuestre que, entonces, tanto
$U$ como $V$ son ideales a izquierda de $A$, es decir,
\[
  A\cdot U = U
  \qquad\text{y}\qquad
  A\cdot V = V.
\]
}

\inlinehead{Respuesta} Probemos primero que $U$ es un ideal a izquierda. Sea
$a\in A$ y $u\in U$. Como $A=\CC[G]$, podemos escribir
\[
  a=\sum_{g\in G}\alpha_g g
  \qquad\text{con } \alpha_g\in\CC.
\]
Entonces
\[
  a\cdot u=\sum_{g\in G}\alpha_g\,(g\cdot u).
\]
Como $U$ es una subrepresentación, se tiene $g\cdot u\in U$ para todo
$g\in G$. Por linealidad, concluimos que $a\cdot u\in U$. Esto prueba que
\[
  A\cdot U\subseteq U.
\]
Por otra parte, como $1\in A$, para todo $u\in U$ se tiene
\[
  u=1\cdot u\in A\cdot U.
\]
Luego $U\subseteq A\cdot U$, y por tanto
\[
  A\cdot U=U.
\]

El mismo argumento aplicado a $V$ muestra que
\[
  A\cdot V=V.
\]
Así, tanto $U$ como $V$ son ideales a izquierda de $A$.\hfill$\blacksquare$

\problem{5}{%
En el contexto del ejercicio anterior, suponga además que $U$ y $V$ no
comparten ``factores irreducibles'' en sus descomposiciones en irreps, es
decir, si
\[
  U = \bigoplus U_i
  \qquad\text{y}\qquad
  V = \bigoplus V_j,
\]
donde todas las $U_i$ y $V_j$ son irreps, entonces se cumple
\[
  U_i \not\cong V_j \qquad \text{para todo } i,j.
\]
Otra forma de decir esto mismo, por el lema de Schur, es que
\[
  \operatorname{Hom}_G(U,V)=0.
\]
Demuestre que entonces $U$ y $V$ también son ideales a derecha, es decir,
\[
  U\cdot A = U
  \qquad\text{y}\qquad
  V\cdot A = V.
\]
}

\inlinehead{Respuesta} Ya sabemos por el ejercicio anterior que $A=\CC[G]$
es una representación de $G$ por multiplicación izquierda, y que
\[
  A=U\oplus V
\]
como suma directa de subrepresentaciones. Denotemos por
\[
  p_U\colon A\to U,
  \qquad
  p_V\colon A\to V
\]
las proyecciones asociadas a esta descomposición; como $U$ y $V$ son
subrepresentaciones, ambas son morfismos de representaciones.

Fijemos $g\in G$ y definamos
\[
  R_g\colon A\to A,
  \qquad
  R_g(x)=xg.
\]
Si $h\in G$ y $x\in A$, entonces
\[
  R_g(h\cdot x)=R_g(hx)=(hx)g=h(xg)=h\cdot R_g(x).
\]
Por lo tanto, $R_g$ es un morfismo de representaciones. Consideremos ahora
\[
  T_g:=p_V\circ R_g|_U\colon U\to V.
\]
Entonces $T_g\in \operatorname{Hom}_G(U,V)$, y por hipótesis
$\operatorname{Hom}_G(U,V)=0$. Luego $T_g=0$. Así, para todo $u\in U$,
\[
  p_V(ug)=0,
\]
y como $A=U\oplus V$, esto implica que $ug\in U$. En particular,
$U\cdot g\subseteq U$ para todo $g\in G$.

Si ahora $a\in A$, escribimos
\[
  a=\sum_{g\in G}\alpha_g g
  \qquad\text{con } \alpha_g\in\CC.
\]
Entonces, para $u\in U$,
\[
  u\cdot a
  =u\left(\sum_{g\in G}\alpha_g g\right)
  =\sum_{g\in G}\alpha_g\,(ug).
\]
Como cada $ug$ pertenece a $U$, por linealidad $u\cdot a\in U$. Por tanto,
\[
  U\cdot A\subseteq U.
\]
La inclusión opuesta es inmediata porque $1\in A$ y, para todo $u\in U$,
\[
  u=u\cdot 1\in U\cdot A.
\]
Así, $U\cdot A=U$.

El mismo argumento aplicado a
\[
  S_g:=p_U\circ R_g|_V\colon V\to U
\]
y usando que también $\operatorname{Hom}_G(V,U)=0$, muestra que $vg\in V$
para todo $v\in V$ y todo $g\in G$. Por linealidad, $V\cdot A\subseteq V$,
y como $v=v\cdot 1$, se obtiene $V\cdot A=V$.

Concluimos que $U$ y $V$ son ideales a derecha de $A$.\hfill$\blacksquare$

\problem{6}{%
(*) Sea $f\colon G\to\CC$ una función de clase. Demuestre, con los detalles,
que
\[
  f * f' = f' * f
\]
para cualquier otra función $f'\colon G\to\CC$.
}

\inlinehead{Respuesta} Usamos la definición usual de convolución:
\[
  (f*f')(g)=\sum_{x\in G} f(x)\,f'(x^{-1}g)
  \qquad\text{para todo } g\in G.
\]
Fijemos $g\in G$. Entonces
\[
  (f*f')(g)=\sum_{x\in G} f(x)\,f'(x^{-1}g).
\]
Hacemos el cambio de variable
\[
  y:=x^{-1}g.
\]
Como la aplicación $x\mapsto x^{-1}g$ es una biyección de $G$ en $G$, podemos
reindexar la suma y obtenemos
\[
  (f*f')(g)=\sum_{y\in G} f(gy^{-1})\,f'(y).
\]
Ahora usamos que $f$ es una función de clase. Observemos que, para todo
$y\in G$,
\[
  gy^{-1}=y\,(y^{-1}g)\,y^{-1},
\]
de modo que $gy^{-1}$ y $y^{-1}g$ son conjugados. Como $f$ es constante en
clases conjugadas, se sigue que
\[
  f(gy^{-1})=f(y^{-1}g).
\]
Sustituyendo esto en la suma anterior,
\[
  (f*f')(g)=\sum_{y\in G} f(y^{-1}g)\,f'(y)
  =\sum_{y\in G} f'(y)\,f(y^{-1}g)
  =(f'*f)(g).
\]
Como esto vale para todo $g\in G$, concluimos que
\[
  f*f'=f'*f.
\]
Esto prueba que toda función de clase conmuta, bajo convolución, con cualquier
otra función $f'\colon G\to\CC$.\hfill$\blacksquare$

\problem{7}{%
Sea $c$ una clase conjugada de $G$, y sea
\[
  s_c = \sum_{g\in c} g \in \CC[G].
\]
Sea $(U,\pi)$ una representación irreducible de $G$. Demuestre que $\pi(s_c)$
es un operador escalar.

\smallskip
\textit{Recuerde que}
\[
  \pi(s_c) = \sum_{g\in c}\pi(g),
\]
por definición.
}

\inlinehead{Respuesta} Sea
\[
  T:=\pi(s_c)=\sum_{g\in c}\pi(g).
\]
Veamos primero que $T$ conmuta con todos los operadores $\pi(h)$, $h\in G$.
En efecto, para todo $h\in G$,
\[
  \pi(h)\,T\,\pi(h)^{-1}
  = \sum_{g\in c}\pi(hgh^{-1}).
\]
Como la aplicación $g\mapsto hgh^{-1}$ es una biyección de $c$ en sí misma,
al reindexar la suma obtenemos
\[
  \sum_{g\in c}\pi(hgh^{-1})=\sum_{g\in c}\pi(g)=T.
\]
Por lo tanto,
\[
  \pi(h)T=T\pi(h)
  \qquad\text{para todo } h\in G.
\]

Ahora repetimos el argumento del Problema~2. Como $U$ es un
$\CC$-espacio vectorial de dimensión finita, $T$ tiene al menos un valor
propio $\lambda\in\CC$. Sea
\[
  W:=\ker(T-\lambda I).
\]
Entonces $W\neq\{0\}$. Además, si $v\in W$ y $h\in G$, entonces
\[
  (T-\lambda I)\bigl(\pi(h)v\bigr)
  =T\pi(h)v-\lambda\pi(h)v
  =\pi(h)Tv-\lambda\pi(h)v
  =\pi(h)(Tv-\lambda v)=0.
\]
Así, $\pi(h)v\in W$, y por tanto $W$ es un subespacio $G$-invariante de $U$.

Como $(U,\pi)$ es irreducible y $W\neq\{0\}$, se sigue que $W=U$. Esto
significa que
\[
  (T-\lambda I)u=0
  \qquad\text{para todo } u\in U,
\]
es decir,
\[
  \pi(s_c)=T=\lambda I.
\]
Concluimos que $\pi(s_c)$ es un operador escalar.\hfill$\blacksquare$

\end{document}
